% Research handoff, prepared with substantial OpenAI GPT-6 Astra/Codex assistance.
% Extends icekylinx PR18/PR24, Zhihao Chen PR21/PR23, Rohan Arun PR25,
% and attributed RaD/hipotures semantic/bulk interfaces. Apache-2.0.
\section{A near-balanced compatible bit basis}
The selected dimensions are
\[
 a=33,\quad b=31,\quad k=a+1=34,\quad H=ab=1023,\quad m=Hk=34782.
\]
Use exactly the retained-total binary scalar construction, positive-label
propagation, carrier matching and dirty-scratch schedule of PR24, at these
three dimensions. The base-four producer is unchanged as an algorithm.
The newly reconstructed local counts $(h,v,c,q,|\mathcal M|,R,\mathrm{loss})$ are
\[
\begin{array}{c|rrrrrr}
33&5456&114729&16401&3054&128076&1056\\
31&4495&93670&13516&2683&104503&930\\
34&5984&127823&17986&3706&142103&1122.
\end{array}
\]
The positive-label matching at $h=34$ has 68 more matches than the initial
envelope matching. This is permissible: the initial continuation arcs are
included before positive enlargement, and every final match separately
passes the exact positive-space inclusion and strict event-order test.
Equality of the two matching cardinalities is unnecessary for the carrier
proof. Each value has an unmatched first use, and no matched continuation
can precede its donor. All saved histograms have nonnegative entries and
\(\sum_r rH_{h,r}=hR_h+2h(h-1)\).

\paragraph{One common rational family.}
Let $I=(0,\ldots,a-1)$ and prescribe
\[
 R=C=[I,0,I],\qquad n=a+k=67.
\]
For inner physical index $i=b\alpha+\beta$, use the same family
$K=T_2(I_a\otimes L_b)$, $J_\beta=M_\beta L_a$ as PR18/24.
Prescribe row $\alpha$ of $M_\beta$ to be coordinate $R_i$ for $i<n$,
and inverse column $\alpha$ to be coordinate $C_j$ at $i=H-n+j$.
The effective residues are $\{0,2,3\}$ on the row side and
$\{26,28,29\}$ on the inverse-column side. They are pairwise distinct
and disjoint. The physical boundary positions are also disjoint.
Every partial prescription completes to a permutation matrix. Thus the
remaining entries form an affine rational space whose invertible locus
is nonempty and irreducible. $L_a,L_b$ remain free.
The first $a$ rows and last $a$ inverse columns have matching $I$ labels;
all ordinary axis-$a$ corners therefore survive. Rank-one inactive-factor
weights separate for the axis-$b$ corners, leaving a diagonal rescaling of
$L_bPL_b^{-1}$, so all ordinary axis-$b$ corners survive as well.

The outer family is $S=T(I_k\otimes K)$ with free $L_k$. Prescribe the
first $k$ outer row functionals in $L_k$ order followed by repetitions of
its first $a$ row functionals. Prescribe the last $k$ inverse columns in
matching $L_k^{-1}$ order, preceded by repetitions of its first $a$
inverse columns. The boundary intervals of length $67$ are disjoint.
Ordinary axis-$k$ corners and the unmatched full $k$-block survive.

\paragraph{The exact A5 incidence criterion.}
The first $a+b-1=63$ inner nullspace rows of
$(I-P_a)\otimes(I-P_b)$ correspond, after nonzero diagonal rescaling,
to edges $(R_i,a+(i\bmod b))$ of a bipartite graph with $a+b=64$
vertices. The last 63 inverse columns correspond to
$(C_j,a+((H-n+j)\bmod b))$, for $n-63\le j<n$.
Both graphs are trees: the accompanying certificate enumerates their
63 edges, checks all vertices and connectivity, and records leaf-removal
orders. A connected graph with 64 vertices and 63 edges is a tree;
leaf elimination gives rank 63 for each restriction. This is an exact
finite combinatorial proof for the selected dimensions, not a floating
rank test. The usual large-projector identity gives the middle block
$H-2(a+b-1)=897$.

\subsection{Two A1 blocks for $k=a+1$}
The paired A1 null projector is
\[
 Q=I_k\otimes P_q+P_t\otimes P_U,\quad
 \operatorname{rank}P_q=1,\quad\operatorname{rank}P_U=a,\quad P_qP_U=0.
\]
Its rank is $n=2a+1=67$. The actual A1 projector is $I-Q$.
Write $P_q=P_{x_a}\otimes P_{x_b}$ and
$P_U=I_a\otimes P_{u_b}$, where $P_{x_b}P_{u_b}=0$.
Absorb the free factor bases into the primal and dual coordinates.
For the first-$n$/last-$n$ corner, the outer row labels are
$p=[I,a,I]$, and the outer inverse-column labels are $q=[I,I,a]$.
After nonzero diagonal row and column rescalings, the corner is exactly
\[
 M_{ij}=\mathbf1_{R_i=C_j}+f_i g_j\mathbf1_{p_i=q_j}.
 \tag{A1}
\]
Indeed the $P_t\otimes P_U$ term supplies the coefficient-label delta,
and the $I_k\otimes P_q$ term supplies the outer-label delta. Each $f_i$
and $g_j$ is a product of a nonzero motif-line coordinate, a ratio of
coordinates of the two middle-axis lines, and an inverse outer-line
coordinate. These are rational functions on the single family above.

Group the rows as $A_0,\ldots,A_{a-1};B;C_0,\ldots,C_{a-1}$,
and the columns as $A'_0,\ldots,A'_{a-1};B';C'_0,\ldots,C'_{a-1}$.
The first $a$ rows against the last $a$ columns give
\[
 T_{ij}=\delta_{ij}+\delta_{i,j+1}f_{A_i}g_{C'_j}.
\]
This is lower bidiagonal with diagonal one. Its pivots are therefore the
contiguous equally ordered block $(i,a+1+i)$, $0\le i<a$.
Left multiplication by its lower-triangular inverse uses only the retained
ordered-affine operations. It realizes one width-$a$ recursive swap.

Here is an explicit Schur computation for the rest. Put
$d_i=1+f_{A_i}g_{A'_i}$ and $v=1+f_{A_0}g_{B'}$.
After eliminating $T$, row $C_i$ has a common factor
$f_{C_i}-f_{A_i}$; divide it out. Its entries on columns $A'_j,B'$ are
\[
 L_{ij}=\begin{cases}
 0&j>i,\\ g_{A'_i}&j=i,\\
 -g_{C'_{i-1}}(T^{-1})_{i-1,j}d_j&j<i,
 \end{cases}
\]
\[
 h_0=g_{B'},\qquad
 h_i=-g_{C'_{i-1}}(T^{-1})_{i-1,0}v\quad(i>0).
\]
Write $z=f_Bg_{C'_{a-1}}$. Row $B$ is
\[
 B_j=-f_{A_0}g_{A'_0}\delta_{j0}
      -z(T^{-1})_{a-1,j}d_j,
\qquad
 w=-f_{A_0}g_{B'}-z(T^{-1})_{a-1,0}v.
\]
The next pivot is $w$ in column $B'$. After eliminating it,
row $C_0$ has rightmost entry
$-g_{B'}B_{a-1}/w$ in column $A'_{a-1}$, giving one further scalar pivot.
For $1\le i\le a-2$, eliminating this latter pivot cancels the rank-one
term exactly and leaves
\[
 L_{ij}-(h_i/g_{B'})L_{0j}.
\]
For $j>i$ this is zero, while at $j=i$ it is $g_{A'_i}\ne0$.
The two preceding scalar pivot columns $B\prime,A\prime_{a-1}$ have
already been eliminated in this Schur complement. Within the new block,
the previously chosen block pivots are at smaller columns $j<i$; those
rows vanish to the right of their pivots and cannot change these entries.
Thus $(C_i,A'_i)$, $1\le i\le a-2$, form one contiguous width-$(a-2)$
block. The last remaining pivot is scalar. The ordered corner profile is
\[
 a,1,1,a-2,1.
\]

\paragraph{Nonvanishing within the actual family.}
The factors $f_{C_i}-f_{A_i}$ compare primal middle-axis coordinate ratios
at shift $k\bmod b=3$, at the same coefficient and outer label. Distinct
coordinate ratios of the two independent middle-axis motif lines are
not identically equal under free $L_b$. For the primal nullspace
restriction, eliminate the first $k$ outer-coordinate rows: the repeated
$a$ rows then give a diagonal matrix with these nonzero differences.
For the dual restriction, eliminate the last $k$ outer-coordinate columns.
The first $a$ repeated columns leave a triangular matrix: coefficient zero
compares dual ratios at shift $a\bmod b=2$, and every other column has a
nonzero diagonal and at most one preceding coefficient. Its determinant
is not identically zero. Hence both nullspace restrictions, and therefore
the complete corner, are generically nonsingular.

The remaining denominators and coordinates are nonzero generically.
The scalar $B_{a-1}=-z d_{a-1}$ is not identically zero. Neither is $w$:
after clearing the outer diagonal weights $\omega_i=t_i\varphi_i$,
its two summands have, respectively, a product of the weights
$\omega_1\cdots\omega_a$ and an affine dependence on $\omega_0$.
These are not identically proportional on $\sum_i\omega_i=1$ for $a>1$.
Thus the displayed scalar pivots are available on a nonempty open set.
The final scalar pivot follows from full nonsingularity. These statements
also have concrete rational witnesses, recorded in full in the certificate.
For $a=33$ and seed 4, the control records every primal/dual vector, all
normalized corner pivot values, the two stated scalar pivots, the last
pivot, and the nonzero corner determinant. The middle-axis vectors are
exactly biorthogonal. Any ordered pair of orthogonal rank-one projectors
is rationally conjugate to these witnesses: map the two primal basis
vectors and their common dual kernel to the target pair and its common
dual kernel. The analogous rank-one conjugacy applies independently to
the first and outer factors. The prescribed $M_\beta$ admit permutation
completions, and the disjoint outer boundary prescriptions admit invertible
rational completions. Thus this is an actual rational point in the same
family, proving nonvanishing of each displayed pivot rational function.
It is not a sample of unrelated corner entries. This explicit witness
can replace the preceding polynomial argument for nonvanishing; it does
not need to satisfy the other, separately witnessed generic minors.

The corner of $I-Q$ differs by a global minus sign, which does not alter
pivot positions. Since it is the full nullity corner, its elimination
leaves the identity middle block $m-2n=34648$. The complete paired A1
profile is therefore
\[
 \boxed{33,1,1,31,1,34648}.
\]
All these are actual consecutive equally ordered pivot intervals; no
uncharged gathering, reversal or independent edge basis is used.

\subsection{The A5 block and the other profiles}
The PR25 A5 observation generalizes directly because the first $a$ row
labels and last $a$ inverse-column labels are still $I$.
In the 63-square A5 null corner, write $p\xi,v\nu$ for the two inner
line projectors. Its exact entries are
\[
 Q_{ij}=\mathbf1_{\beta_i=\beta'_j}p_{r_i}\xi_{c_j}
 +\mathbf1_{r_i=c_j}v_{\beta_i}\nu_{\beta'_j}
 -p_{r_i}\xi_{c_j}v_{\beta_i}\nu_{\beta'_j}.
\]
The first pivot $(0,62)$ is pure rank one. After its elimination, rows
$1\le i\le29$ have rightmost pivots at column $i+32$, equal to
$p_i\xi_{i+2}\ne0$. The possible coefficient-label term is two columns
to their left. This proves a width-29 block; the other 34 corner pivots
may remain scalar. Together with its middle block, A5 has profile
$34$ singletons and widths $29,897$. This contribution is attributed to
Rohan Arun's PR25; only its changed dimension is specialized here.

The remaining PR24 corner arguments use only matched outer intervals,
nonzero line coordinates, and free second rows/next-to-last inverse columns.
They apply to the present same irreducible family: the first/last $k$
outer intervals are disjoint, $H>2k$, and $H>2(a+k)$. In particular the
A3 full nullspace restriction is witnessed by those free second rows;
its inner corner is a rescaled rank-$(k-1)$ projector, with one scalar
and a $(k-2)$ block. After that perturbation, the matching interval has
width $H-2k+2$. Its remaining $2(k-1)$ corner pivots remain scalar.
The large-projector identity leaves $m-2H-2k+2$.
The gauged stage-two null corner is diagonal of width $b$, by the same
factorization as PR21/24. All its factors are nonzero coordinates in
the prescribed family. Thus the complete macro/physical table is
\[
\begin{array}{c|c|l}
\text{class}&\text{copies}&\text{widths}\\\hline
A1&B_1&33,1,1,31,1,34648\\
\text{unmatched stage three}&B_3-B_1&34,34714\\
\text{gauged stage two}&B_2&31,34720\\
A3&2N&67\text{ singletons};32,957,32670\\
A5&2N&34\text{ singletons};29,897\\
\text{data growth}&2N\text{ each}&(1,31),(1,29),(1,32).
\end{array}
\]
For each axis $h$, replicate histogram entry $H_{h,r}$ by $N/v_h$.
If $2r>h$, use $h-r$ scalar children and one width-$2r-h$ child;
otherwise use $r$ scalar children. This accounts for every edge.
Each required generic condition is a nonzero rational function on the
same irreducible parameter family. Their finite product is nonzero.
A simultaneous rational point can therefore be found by enumeration;
then choose an odd address prime avoiding its finitely many bad minors
and denominators. The construction is independent of recursive width.

\subsection{Physical schedule and analytic composition}
Set $N=v_av_bv_k$ and $B_i=NR_{h_i}/v_{h_i}$.
The rebuilt counts satisfy $B_3\ge B_1$. The same regular intersection-one
neighbor graph provides a perfect matching of middle-axis triples,
so the stage-one bank shares with a subset of stage three. The unchanged
common-frame dirty-scratch word restores every auxiliary value.
Use the PR21/24 source gauge $D_{P_0}$ and sink $D_{I-P_0}$ on every
stage-two auxiliary: its entrance is identity and its exit rank is $m-b$.
The physical data enlargement precedes the first inverse copy as in PR24.
The rank identity is
\[
 W=2N+B_2+B_3=7190435649232,\qquad
 s=Wm-2N+2L,\quad L=N\sum_{h\in\{33,31,34\}}6/(h-2),
\]
\[
 Wm-s=120943042880.
\]
The common basis above compiles the exact physical profiles; spectators,
integer-width floors, bounded tails, arbitrary dirty scratch, one-time
row padding and fixed-tape scheduling follow the retained arbitrary-width
bit interchange proof. Every child is proper and the maximum is 34720.

The independent complex layer remains exactly PR21's $h=28$ circuit,
with $a_c=18/10^6$, $m_c=21952$, $W_c=2085111546336$, and maximum child
21924. Its complete physical/phase evidence and semantic guard $C_1=1$
are unchanged. Only the bit stock changes in PR23's product row reservoir.
The certificate recomputes the bit halving depth, bit wire exponent,
product-stock degree, and compressed eventual checks. It then uses the
PR23 semantic/bulk assembly with $a_b=52007/4000000000$, $\beta=1/4$,
$h=10^{-8}$, $q=a_b(1-2h)$, $c=q(1+h)$,
$\epsilon=(1-h)/(1+c+q)$, $\lambda'=1-q$,
$\lambda=(1-a_b+\lambda')/2$, and
$r=(\epsilon q+1-\epsilon)/2$, $\delta=h/8$.
All 47 strict constraints and seven margins are recomputed exactly.
The chosen final witness is
\[
 \boxed{\kappa=13001411/10^{12}=1.3001411\times10^{-5}}.
\]
This assumes the same upstream multiplication theorem, ordered-affine
compiler, arbitrary-width/common-basis transfer, fixed finite-alphabet tape
scheduler, and the attributed RaD/PR20/PR23 semantic, routing, phase-cell
inverse, bulk locality, prime-supply and recovery proofs. It changes the
bit construction; it does not improve the complex circuit or introduce
new analytic assumptions beyond those explicitly inherited from PR23.
Exact arithmetic, finite graph checks, and the written generic-basis proof
do not constitute formal verification of the complete multiplication theorem.
